%! Advanced Factoring — Unit 3, Lesson 3.2 — Homework
\documentclass[10pt]{article}
\usepackage{algebra2-article}
\usepackage{algebra2-boxes}
\newcommand{\TallMath}[1]{$\displaystyle #1\rule[-1.4em]{0pt}{3.2em}$}

\begin{document}

\pageheader{Unit 3, Lesson 3.2}{Homework}
\namedateperiod

\vspace{0.10in}

\begin{notesbox}{Practice}
Factor \textbf{completely} every time. Take the GCF out first, then count the terms to choose your
tool --- and check each factor before you call a problem done.

\begin{enumerate}[leftmargin=1.9em, itemsep=11pt]
  \item \textbf{GCF first.}
    \begin{enumerate}[label=(\alph*), leftmargin=1.8em, itemsep=5pt]
      \item $18x^4 - 24x^3 + 30x^2 = $ \blank{5.5cm}
      \item $14x^3y - 21x^2y^2 = $ \blank{5.5cm}
      \item Factor out $-5x$: \; $-5x^3 + 20x^2 - 35x = $ \blank{5.0cm}
    \end{enumerate}

  \item \textbf{Two terms.} Name the pattern, then factor.
    \begin{enumerate}[label=(\alph*), leftmargin=1.8em, itemsep=5pt]
      \item $25x^2 - 81 = $ \blank{4.6cm}
      \item $x^3 + 64 = $ \blank{4.6cm}
      \item $8x^3 - 125y^3 = $ \blank{4.6cm}
      \item $x^2 + 49 = $ \blank{4.6cm}
    \end{enumerate}

  \item \textbf{Three terms.}
    \begin{enumerate}[label=(\alph*), leftmargin=1.8em, itemsep=5pt]
      \item $x^2 + 16x + 64 = $ \blank{4.6cm}
      \item $x^2 - 2xy - 15y^2 = $ \blank{4.6cm}
      \item $2x^2 + 11x + 12 = $ \blank{4.6cm} \quad (grouping: product $24$, sum $11$)
    \end{enumerate}

  \item \textbf{Four terms --- grouping.}
    \begin{enumerate}[label=(\alph*), leftmargin=1.8em, itemsep=5pt]
      \item $x^3 + 3x^2 + 7x + 21 = $ \blank{5.0cm}
      \item $x^3 - 4x^2 - 9x + 36 = $ \blank{5.0cm} \quad (this one keeps going!)
    \end{enumerate}

  \item \textbf{Factor completely --- more than one step.}
    \begin{enumerate}[label=(\alph*), leftmargin=1.8em, itemsep=5pt]
      \item $2x^4 - 32 = $ \blank{5.5cm}
      \item $x^4 - 10x^2 + 9 = $ \blank{5.5cm}
    \end{enumerate}

  \item \textbf{Back to a graph you already read.} In Lesson 3.0 you read the graph of
        $q(x) = x^4 - 5x^2 + 4$ and found four $x$-intercepts.
    \begin{enumerate}[label=(\alph*), leftmargin=1.8em, itemsep=5pt]
      \item Factor $q(x)$ completely: \blank{5.5cm}
      \item List the zeros: \blank{4.4cm}
      \item Substitute $x = 0$ into your factored form to find the $y$-intercept: \blank{2.0cm}
            Does it agree with the graph? \blank{1.4cm}
    \end{enumerate}

  \item \textbf{Explain.} $x^2 + 9$ cannot be factored, but $x^3 + 27$ can. Explain the difference,
        and include the factorization of $x^3+27$ as part of your answer. \writelines{3}
\end{enumerate}
\end{notesbox}

\begin{extensionbox}
\textbf{Prove it, then use it.}
\begin{enumerate}[label=(\alph*), leftmargin=1.8em, itemsep=5pt]
  \item Multiply $(a + b)(a^2 - ab + b^2)$ out completely --- all six products --- and show which four
        cancel. What is left?
  \item A solid cube of ice measures $x$ cm on each edge. A cubical hole $3$ cm on each edge is
        carved out of one corner. Write the remaining volume as a difference of two cubes, then
        factor it.
  \item Evaluate your factored expression at $x = 5$ cm, and check it against
        $5^3 - 3^3$. Do they agree?
  \item The factor $(x - 3)$ is the difference of the two edge lengths. Explain why it makes sense
        that the volume must contain that factor --- what happens to the leftover ice when $x = 3$?
\end{enumerate}
\writelines{4}
\end{extensionbox}

\begin{spiralbox}
\textbf{Coming next --- Lesson 3.3: Dividing Polynomials.} Today's toolkit is powerful but it has a
boundary: nothing you learned cracks $x^3 - 3x + 2$, even though you know from Lesson 3.0 that it
equals $(x-1)^2(x+2)$. Next you will learn to \emph{divide} one polynomial by another --- long
division and the faster synthetic division --- and the \textbf{Remainder} and \textbf{Factor
Theorems}, which let you test whether $(x-r)$ is a factor before you divide. That is the tool that
gets past today's wall.
\end{spiralbox}

\end{document}
